\cvname{Yan Lin}

\cvtagline{Assistant Professor / Department of Computer Science / Aalborg University}

Email: \href{mailto:lyan@cs.aau.dk}{lyan@cs.aau.dk} \,/\,
Homepage: \href{https://www.yanlincs.com}{yanlincs.com} \,/\,
Phone: +45 9186-1833

{\color{rule}\rule{\linewidth}{0.4pt}}

\vspace{\entryskip}
I am a researcher in computer science.
My research centers on data mining, i.e., the extraction and use of patterns from large-scale data with machine learning models.
My current focus spans trajectory representation learning, spatiotemporal data mining, and AI for materials science.
I am building my teaching experience, drawing on both my research background and my hobbyist exploration of computer science.

\vspace{\entryskip}
I have authored 33 publications in top-tier venues such as SIGMOD, KDD, NeurIPS, AAAI, and IEEE TKDE.
Of these, 17 are first, co-first, or lead student author papers.
My publications have over 1,000 citations on \href{https://scholar.google.com/citations?user=nHMmG2UAAAAJ}{Google Scholar}.
My h-index is 15 and i10-index is 23.


\section*{Education}

\roleentry{Beijing Jiaotong University}{Ph.D. in Computer Science}{Sep 2019 -- Mar 2024}
{Thesis: Self-supervised Learning of Spatiotemporal Trajectory Data. Advisor: Prof. Huaiyu Wan.}

% \roleentry{Aalborg University}{Visiting Ph.D. in Computer Science}{Sep 2022 -- Sep 2023}
% {Advisors: Prof. Bin Yang and Prof. Jilin Hu.}

\roleentry{Beijing Jiaotong University}{B.Sc. in Computer Science}{Sep 2015 -- Jun 2019}{}


\section*{Experience}

\roleentry{Department of Computer Science, Aalborg University}{Assistant Professor}{Jun 2026 -- present}
{}

\roleentry{Department of Computer Science, Aalborg University}{Postdoctoral Researcher}{Jun 2024 -- May 2026}
{Worked on AI for materials science, in collaboration with the AAU Glass Structure and Mechanics group.}
